\documentclass[10pt,a4paper]{amsart}
\usepackage[margin=19mm]{geometry}
\usepackage{amsmath,amssymb,mathtools}
\usepackage[scr=rsfs,cal=euler]{mathalfa}
\usepackage{xfrac}
\usepackage[unicode,hidelinks,bookmarks=false]{hyperref}
\newcommand{\ie}{i.e.\ }
\newcommand{\cf}{cf.\ }
\newcommand{\resp}{respectively\ }
\newcommand{\Subsection}{\subsection}
\providecommand{\nobreakdash}{\nobreak\hbox{-}}
\setlength{\emergencystretch}{2em}
\multlinegap0pt
\usepackage{amssymb}
\usepackage{dsfont}
\newtheorem{thm}{Theorem}
\newtheorem{cor}{Corollary}
\title{Evaluating time-to-event forecasts under fixed-horizon right-censoring}
\author{Robert J. Taggart, Nicholas Loveday, Simon Louis}
\date{Source: arXiv:2603.14835v2 (CC BY 4.0)}
\begin{document}
\maketitle

\noindent\emph{Source and licence.} Evaluating time-to-event forecasts under fixed-horizon right-censoring. Version arXiv:2603.14835v2 (CC BY 4.0), distributed by arXiv under the Creative Commons Attribution 4.0 International licence, http://creativecommons.org/licenses/by/4.0/, as recorded in the arXiv licence metadata for this exact version and shown on the abstract page https://arxiv.org/abs/2603.14835v2. Full licence notice: \texttt{licenses/arxiv-2603.14835-CC-BY-4.0.txt}.\par\smallskip

\noindent\textit{Editorial scope.} Counted unit: the article's cor elicitable (source0.tex lines 589--591). The article's complete original proof of it is delivered with it (source0.tex lines 860--877, one \texttt{proof} environment of 18 lines, which the article places away from the statement). No mathematics is written, repaired or re-derived: the statement and the proof are the article's own lines, in the article's own notation, and no retained line is substituted or re-flowed.  Retained uncounted as the source-local context the proof needs: lines 579--587.  Nothing was omitted from any retained span, so the package carries no omission record.  No retained line contains a citation, so the wrapper carries no bibliography and no bibliography entry.  No retained line contains a figure environment, an image-inclusion command or drawing code, so no image is delivered and the package declares no asset dependency.  Editorial formatting only, in addition: an AMS \texttt{amsart} preamble that restates the article's own declarations and macros used by the retained lines, namely the environments \texttt{thm}, \texttt{cor} on separate counters, together with the standard packages those lines need.  The retained environments print in this document as Theorem 1, Corollary 1 in order of appearance, whereas the article prints the same items with the numbering of its own full document.  The complete native source of the article as distributed in the arXiv e-print for this exact version is bundled unchanged as \texttt{sources/196436/source0.tex}.
\begin{thm}\label{thm:not prov elic}
Suppose that $\tau$ belongs to the interior of $I$, $\mathcal{F}$ is a class of probability distributions on $I$ and $U:\mathcal{F}\to\mathcal{P}(I)$ is a functional on $\mathcal{F}$. If $\mathcal{F}$ contains two distributions $F_1$ and $F_2$ such that
\begin{enumerate}
\item[(a)] $U(F_1) \subset(0,\tau)$ and $U(F_2) \subset(0,\tau)$ 
\item[(b)] $U(F_1) \neq U(F_2)$ and
\item[(c)] $[F_1]_\tau = [F_2]_\tau$,
\end{enumerate}
then $U$ is not locally elicitable with respect to $\mathcal{F}$ on $[0,\tau)$.
\end{thm}

\begin{cor}\label{cor:expectation not provisionally elicitable}
Suppose that $\tau$ is in the interior of $I$ and that $m$ is a positive integer. The mean functional $U$ is not locally elicitable on $[0,\tau)$ relative to any class $\mathcal{F}$ of probability distributions on $I$ that contains either (i) the distributions on $I$ with finite support, or (ii) the distributions on $I$ with a probability density function and finite $m$th moment.
\end{cor}

\begin{proof}[Proof of Corollary~\ref{cor:expectation not provisionally elicitable}]
Denote by $U$ the expectation functional and suppose that $\tau$ lies in the interior of $I$.
Choose $a$ in the interior of $I$ such that $\tau<a<2\tau$. Let $f_1$ and $f_2$ denote the probability densities given by
\begin{align*}
f_1(s)&=
\begin{cases}
1/a & \text{if } 0\leq s \leq a \\
0 & \text{otherwise},
\end{cases} \\
f_2(s)&=
\begin{cases}
1/a & \text{if } 0\leq s \leq \tau \\
\frac{2}{a}(1 - (s-\tau)/(a-\tau)) & \text{if } \tau < s \leq a \\
0 & \text{otherwise,}
\end{cases}
\end{align*}
and let $F_1$ and $F_2$ denote their respective CDFs. Note that $F_1$ and $F_2$ have finite support, density functions and finite $m$th moment. Now $0 < U(F_2) < U(F_1)=a/2<\tau$, whence $F_1$ and $F_2$ satisfy properties (a), (b) and (c) of Theorem~\ref{thm:not prov elic}.  Hence $U$ is not locally elicitable with respect to $\mathcal{F}$ on $[0,\tau)$.
\end{proof}

\end{document}
